\section{Tractability}
\label{sec:tractability}

The foundational paper establishes that exact $\vol(G)$ computation is
\#P-hard on continuous capability spaces \citep{dyer1988complexity}. The
lattice construction replaces this with a finite, structured computation.

\begin{proposition}[Polynomial-Time Computation of $\volL$]
\label{prop:tractability}
For a finite capability lattice $(\C, \leq)$ with $|\C| = n$ nodes,
    $\volL(G)$ is computable in $O(n^3)$ time via the standard M\"{o}bius
    recurrence, or $O(n^{2.37})$ via matrix methods
    \citep{stanley1997enumerative}. For tree-structured lattices (no node has
    multiple covers), the cost reduces to $O(n \cdot k)$ where $k$ is the
    maximum in-degree.
\end{proposition}

\begin{proof}
The computation proceeds in three passes (full algorithm in
Appendix~\ref{app:proof_tractability}):

\emph{Pass 1: Compute the M\"{o}bius function.} Process the lattice in
topological order. For each pair $d' \leq d$, compute $\mu(d', d)$ recursively:
$\mu(d, d) = 1$ and $\mu(d', d) = -\sum_{d' \leq d'' < d} \mu(d', d'')$. There
are $O(n^2)$ comparable pairs, and each inner sum iterates over up to $O(n)$
intermediate elements, giving $O(n^3)$ total. On a tree (each node has at most
one cover), $\mu$ is nonzero only between covers, reducing to $O(n \cdot k)$.

\emph{Pass 2: Compute independent weights.} For each node $d$:
$\iw(d) = \sum_{d' \leq d} \mu(d', d) \cdot w(d')$. Each sum has at most $n$
terms. Total: $O(n^2)$.

\emph{Pass 3: Sum over downward closure.} Given $G \subseteq \C$, compute
$\dc G$ by traversing downward from each element of $G$. Then $\volL(G) =
\sum_{d \in \dc G} \iw(d)$. Cost: $O(n)$.

The total is $O(n^3)$, dominated by Pass~1. For tree-structured lattices,
the total is $O(n \cdot k)$. In practice, capability lattices are sparse
(most capabilities have few comparable elements), and the effective cost is
much closer to $O(n^2)$ than $O(n^3)$.
\end{proof}

\begin{corollary}
\label{cor:sign_tractable}
The sign of $\Delta \volL(G \mid \pi)$ for a proposed action $\pi$ is
    computable in $O(n^3)$ time ($O(n \cdot k)$ for trees), making the full
    optimization loop of the foundational paper (Section~4.2 of
    \citet{lasser2026gfm}) tractable on the lattice.
\end{corollary}

\paragraph{Contrast with \#P-hardness.} The continuous volume problem is hard
because it requires integrating over an arbitrary region of high-dimensional
space. The lattice construction avoids this by discretizing the capability
space: $\volL$ is a finite sum, not an integral. The tradeoff is
expressiveness: the lattice cannot represent arbitrary capability regions, only
collections of discrete, benchmarked capabilities. This is the right tradeoff
for an implementable system, since the GFM actor can only reason about
capabilities it has learned about through communication and observation.
